\section{Introduction}\label{sec:intro}

% Para 1 -- why LV FSI
Computational models of the heart are moving from research tools towards
clinical and regulatory use \citep{Quarteroni2017,Chabiniok2016}. Within these,
fluid--structure interaction (FSI) models of the left ventricle (LV) resolve
the two-way coupling between the contracting myocardium and intraventricular
blood flow, and are used to study diastolic vortex formation, intraventricular
pressure differences, washout and thrombogenic risk, and the haemodynamic
consequences of altered wall mechanics \citep{Mittal2016}.
\TODO{3--5 representative LV FSI references: ALE (e.g.\ Bucelli/Zingaro et al.),
immersed boundary (e.g.\ Gao, Feng, Griffith), immersed/fictitious domain,
FV/OpenFOAM-based.}

% Para 2 -- reproducibility problem
The credibility of such models rests on verification: demonstrating that the
discrete solution converges to the solution of the stated mathematical model
\citep{Roache1998,Oberkampf2010}. In LV FSI this is rarely shown. A survey of
the literature (\cref{tab:survey}) shows that \TODO{X of Y} studies do not
report the ventricular dimensions, \TODO{X of Y} report no mesh-sensitivity
study, \TODO{X of Y} report no time-step-sensitivity study, \TODO{X of Y}
do not state the spatial/temporal discretisation schemes, and \TODO{X of Y}
do not state the stabilisation, the coupling algorithm or its tolerances. Results can therefore be neither reproduced nor compared across
codes.

% Para 3 -- benchmarks that exist
Benchmarks have played this role elsewhere. For general FSI, the
\citet{TurekHron2006} channel-flow benchmark and the experimental benchmark of
\citet{Hessenthaler2017} are widely used, but neither involves an active,
anisotropic, nearly incompressible solid enclosing a cavity. For cardiac solid
mechanics, \citet{Land2015} proposed passive inflation and active contraction
of an idealised LV, and \citet{Arostica2025} extended this to elastodynamics
with a visco-active Holzapfel--Ogden material, analytical fibres and
pericardial boundary conditions, with reference solutions from nine groups.
No equivalent exists for ventricular FSI.
\TODO{Check and position against: experimental LV FSI benchmarks (e.g.\ the
2026 Phys.\ Fluids immersed-particle benchmark, Unipa), Domenichini-type
prescribed-motion LV CFD, and any recent LV FSI code-comparison.}

% Para 4 -- difficulties specific to cardiac FSI that a benchmark must exercise
Cardiac FSI is numerically demanding: the solid is nearly incompressible,
anisotropic and active; collocated discretisations of such a solid need
stabilisation whose magnitude can stiffen the response; blood and myocardium
have similar densities, so partitioned schemes suffer from the added-mass
effect \citep{Causin2005}; and when the cavity is closed by kinematic
(flow-rate) conditions, the enclosed incompressible fluid gives rise to the
incompressibility dilemma of Dirichlet--Neumann coupling \citep{Kuttler2006}.
The settings chosen to deal with these issues --- stabilisation factors,
interface conditions, relaxation and tolerances --- change the answer, yet
they are almost never reported.

% Para 5 -- contribution
This paper proposes a valve-free FSI benchmark for an idealised LV that adds
blood flow to Benchmark~1 of \citet{Arostica2025}. Rather than a full cardiac
cycle, it poses passive filling and active contraction as two separate
problems, each the FSI counterpart of a solid-only AR25 case. Its
contributions are:
\begin{enumerate}
\item a complete, unambiguous problem definition --- geometry and port
  dimensions, fibres, solid and fluid models, interface and boundary
  conditions --- that requires no valve model (\cref{sec:problem});
\item two FSI problems, passive filling and active contraction, each with a
  solid-only counterpart solved by nine codes in AR25, so that the
  effect of the fluid is isolated and verifiable (\cref{sec:steps});
\item precisely defined quantities of interest and a discrepancy metric
  compatible with \citet{Arostica2025} (\cref{sec:qoi});
\item a complete record of the numerical setup --- solid stabilisation,
  fluid schemes, mesh motion, interface conditions, coupling acceleration
  and tolerances (\cref{sec:methods}) --- and systematic studies of mesh,
  time step and each of these settings, giving grid-convergence indices,
  a robustness map and a recommended configuration
  (\cref{sec:verification,sec:results});
\item open meshes, case files and reference data, and a minimum reporting
  checklist for ventricular FSI (\cref{app:checklist}).
\end{enumerate}

The remainder of the paper is organised as follows. \TODO{roadmap}.

% ---------------------------------------------------------------------------
\begin{table}[t]
\centering
\caption{Reporting practice in left-ventricular FSI studies.
\checkmark: reported; --: not reported; (\checkmark): partially.
\TODO{Systematic search: Scopus/WoS, ``left ventricle'' AND ``fluid-structure
interaction'', 2010--2026; record protocol and inclusion criteria.}}
\label{tab:survey}
\small
\begin{tabular}{lccccccc}
\toprule
Study & Method & LV dims & Mesh & $\dt$ & Schemes & Stab. & Coupling \\
\midrule
\TODO{ref} & ALE / IB / FD & & & & & & \\
\TODO{ref} & & & & & & & \\
\midrule
This work & ALE-FV & \checkmark & \checkmark & \checkmark & \checkmark & \checkmark & \checkmark \\
\bottomrule
\end{tabular}
\end{table}
